% Originally: content/week-06.tex, \section{The implicit function theorem} (Corsin Week 6)

% Source: Corsin Nick/Class Notes/Week 6.pdf, pp. 10--12
\section{The implicit function theorem}
\label{sec:implicit_function_theorem}

\begin{theorem}[Implicit function theorem]
\label{thm:implicit_function_theorem}
Let $1 \leq n-d < n$, let $f \in C^k(\mathbb{R}^d\times\mathbb{R}^{n-d}, \mathbb{R}^{n-d})$ (or an
open subset of $\mathbb{R}^d\times\mathbb{R}^{n-d}$), and suppose $f(x_0,y_0) = 0$ for some
$(x_0,y_0) \in \mathbb{R}^d\times\mathbb{R}^{n-d}$ such that the differential $D_yf(x_0,y_0)$ of
$y \mapsto f(x_0,y)$ at $y_0$ is bijective (so $\Jac_yf(x_0,y_0)$ is invertible). Then there exist
balls $B_r(x_0) \subseteq \mathbb{R}^d$ and $B_s(y_0) \subseteq \mathbb{R}^{n-d}$ and
$g \in C^k(B_r(x_0), B_s(y_0))$ such that, for all $(x,y) \in B_r(x_0)\times B_s(y_0)$,
\[f(x,y) = 0 \iff y = g(x).\]
Furthermore,
\[\Jac g(x) = -\big(\Jac_yf(x,g(x))\big)^{-1}\Jac_xf(x,g(x)).\]
\end{theorem}

% Sonnet 5 (Medium)
In words: the equation $f(x,y)=0$ ties $x$ and $y$ together without telling us what $y$ is. The
theorem says that near a point where the tie is \emph{non-degenerate in $y$}, the equation can be
untangled --- there is an honest function $g$ producing the right $y$ for each nearby $x$, it is
as smooth as $f$ was, and its derivative is computable without ever writing $g$ down.

Three things are worth reading off the statement before the pictures.
\begin{itemize}
  \item \textbf{Which variables get solved for is a choice.} Splitting the $n$ coordinates into
    $d$ \qt{free} ones and $n-d$ \qt{solved} ones is ours to make; the theorem only demands that
    the block $\Jac_yf$ belonging to the solved variables be invertible. Choose a different
    split and a different block has to be invertible.
  \item \textbf{The conclusion is an equivalence, not just an inclusion.} Inside the box,
    $f(x,y)=0 \iff y=g(x)$. So $g$ does not merely produce \emph{some} solutions: within the
    box there are no others.
  \item \textbf{Nothing is claimed globally.} The level set $\{f(x,y)=0\}$ is a graph over $x$
    only near the chosen point --- \cref{sec:submanifolds} will call such a set a
    submanifold, and this local-graph property is exactly what that definition abstracts.
\end{itemize}

% Supplement: Sascha Brack/Class Notes/Week_06_Notes_Friday.pdf, pp. 3--5
\begin{example}[A concrete example in $\mathbb{R}^2$]
\label{ex:implicit_function_circle}
Consider $f : \mathbb{R}^2 \to \mathbb{R}$ defined by $f(x,y) := x^2 + y^2 - 1$.
Here $n=2$ and $d=1$, so $x \in \mathbb{R}$ and $y \in \mathbb{R}$. We want to see whether we can write $y$ locally as a function of $x$, i.e., find $g(x)$ such that $f(x,g(x)) = 0$.

First, let us find a point $(x_0,y_0)$ where $f(x_0,y_0) = 0$. We guess $f(0,1) = 0^2+1^2-1 = 0$, so $(0,1)$ works.
Next, we check whether the $y$-derivative $\Jac_y f(0,1)$ is invertible. Since $y$ is 1-dimensional, this is just the partial derivative $\partial_y f(0,1)$:
\[\partial_y f(0,1) = \partial_y(x^2+y^2-1)\big|_{(0,1)} = 2y\big|_{y=1} = 2 \neq 0.\]
Since it is non-zero, it is invertible as a $1\times 1$ matrix.
The theorem then guarantees that for small $r, s > 0$, there exists $g : B_r(0) \to B_s(1)$ such that for all $(x,y) \in B_r(0)\times B_s(1)$,
\[f(x,y) = 0 \iff y = g(x).\]
Since $f \in C^\infty(\mathbb{R}^2)$, we also have $g \in C^\infty(B_r(0))$.
Moreover, we can compute $g'(x) = \Jac g(x)$ using the formula:
\[g'(x) = - \big(\partial_y f(x,g(x))\big)^{-1} \partial_x f(x,g(x)) = - (2g(x))^{-1} (2x) = -\frac{x}{g(x)}.\]
(Indeed, since $y=g(x) = \sqrt{1-x^2}$ near $(0,1)$, its derivative is exactly $-x/\sqrt{1-x^2}$.)
\end{example}

\begin{remark}[Why $\Jac g$ has this form]
The formula is just implicit differentiation, done in blocks. By construction $g$ satisfies
$f(x,g(x)) \equiv 0$ for every $x \in B_r(x_0)$: the left-hand side is identically zero as a
function of $x$, so its derivative is too. Differentiating with the chain rule
(\cref{thm:chain_rule}), and remembering that $x$ enters both directly and through $g$,
\[0 = \underbrace{D_xf(x,g(x))}_{\text{direct}} + \underbrace{D_yf(x,g(x))\cdot Dg(x)}_{\text{through } g},\]
and since $D_yf(x,g(x))$ is invertible near $(x_0,y_0)$ we may move it to the other side:
\[Dg(x) = -\big(D_yf(x,g(x))\big)^{-1} D_xf(x,g(x)),\]
which is the formula in \cref{thm:implicit_function_theorem}.

Read it as an exchange rate. Nudging $x$ changes $f$ by $D_xf$; to keep the equation satisfied,
$y$ must move by whatever compensates, and $\big(D_yf\big)^{-1}$ converts \qt{how much $f$ moved}
back into \qt{how far $y$ must travel}. The minus sign is the compensation. This is also why
invertibility of $D_yf$ is the hypothesis: if $D_yf$ were singular, some change in $f$ could not
be absorbed by moving $y$ at all, and there would be nothing for $g$ to do.
\end{remark}

\begin{ainote}
Corsin remarks in his own notes that ``the geometrical meaning of this one eludes me, sorry''.
The exchange-rate reading above and the two figures below are this document's attempt at the
geometry he was after; the algebra itself is his chain-rule derivation.
\end{ainote}

% Sonnet 5 (Medium) --- FIG-W06-04, replacing the original arrow diagram.
% The hypothesis "D_y f invertible" is a statement about the DIRECTION of the level set at
% the point, so the figure shows the level set and its tangent, not an abstract arrow.
\subsubsection*{What the hypothesis means}

The condition on $D_yf(x_0,y_0)$ is a statement about the \emph{direction} in which the level
set $\{f=0\}$ runs at the point. Moving off $(x_0,y_0)$ straight up, in the $y$-direction alone,
changes $f$ to first order by $D_yf(x_0,y_0)$. If that map is invertible, no vertical motion
leaves $f$ unchanged, so the level set cannot run vertically at $(x_0,y_0)$ --- and a curve that
is not vertical is locally a graph over $x$. If instead $D_yf(x_0,y_0)=0$, the level set is free
to turn vertical, and above a single $x$ there may sit two values of $y$, or none.

\begin{center}
\begin{tikzpicture}[>=Latex, scale=1.05]
  % --- Left panel: D_y f invertible, level set transverse to the vertical ---
  \begin{scope}
    \draw[green!60!black, thick, ->] (-0.3,0) -- (3.5,0) node[right, font=\small] {$x$};
    \draw[orange!90!black, thick, ->] (0,-0.3) -- (0,3.2) node[above, font=\small] {$y$};

    \draw[blue!80!black, very thick, domain=0.35:3.1, samples=60, smooth]
      plot ({\x}, {1.55 + 0.55*(\x-1.6) - 0.16*(\x-1.6)*(\x-1.6)});
    \node[blue!80!black, font=\scriptsize, below right] at (2.75,1.95) {$\{f=0\}$};

    \coordinate (P) at (1.6,1.55);
    \draw[dotted, gray] (1.6,0) -- (P) -- (0,1.55);
    \node[below, font=\scriptsize] at (1.6,-0.05) {$x_0$};
    \node[left, font=\scriptsize] at (-0.05,1.55) {$y_0$};
    \fill (P) circle (1.7pt);

    % vertical probe: moving in y alone leaves the level set
    \draw[red!80!black, very thick, ->] (P) -- ++(0,0.85);
    \node[red!80!black, font=\scriptsize, align=center] at (1.6,2.78)
      {move in $y$ only:\\ $f$ changes};

    \node[font=\scriptsize, align=center] at (1.55,-0.85)
      {$D_yf(x_0,y_0)$ invertible\\ \textbf{solvable for $y$}};
  \end{scope}

  % --- Right panel: D_y f = 0, level set vertical, no graph over x ---
  \begin{scope}[shift={(5.4,0)}]
    \draw[green!60!black, thick, ->] (-0.3,0) -- (3.5,0) node[right, font=\small] {$x$};
    \draw[orange!90!black, thick, ->] (0,-0.3) -- (0,3.2) node[above, font=\small] {$y$};

    \draw[blue!80!black, very thick, domain=0.35:2.75, samples=60, smooth]
      plot ({1.6 - 1.05*(\x-1.55)*(\x-1.55)}, {\x});
    \node[blue!80!black, font=\scriptsize, right] at (0.30,0.38) {$\{f=0\}$};

    \coordinate (Q) at (1.6,1.55);
    \draw[dotted, gray] (1.6,0) -- (Q) -- (0,1.55);
    \node[below, font=\scriptsize] at (1.6,-0.05) {$x_0$};
    \node[left, font=\scriptsize] at (-0.05,1.55) {$y_0$};
    \fill (Q) circle (1.7pt);

    \draw[red!80!black, very thick, ->] (Q) -- ++(0,0.85);
    \node[red!80!black, font=\scriptsize, align=center] at (1.75,2.95)
      {move in $y$ only: $f$ unchanged to first order};

    % two y's over one x, just left of x_0; none just right of it
    \draw[gray, dashed] (1.15,0) -- (1.15,2.75);
    \fill[orange!90!black] (1.15,0.89) circle (1.6pt);
    \fill[orange!90!black] (1.15,2.21) circle (1.6pt);
    \node[orange!90!black, font=\scriptsize, align=center] at (0.62,1.55)
      {two $y$'s\\ here};
    \draw[gray, dashed] (2.15,0) -- (2.15,2.75);
    \node[orange!90!black, font=\scriptsize, align=center] at (2.75,1.15)
      {none\\ here};

    \node[font=\scriptsize, align=center] at (1.55,-0.85)
      {$D_yf(x_0,y_0)=0$\\ \textbf{not solvable for $y$}};
  \end{scope}
\end{tikzpicture}
\end{center}

% Sonnet 5 (Medium) --- FIG-W06-05, replacing the original box-and-arrow diagram.
\subsubsection*{What the conclusion gives}

The theorem does not say the whole level set is a graph --- only that it is one \emph{inside a
small enough box}. On $B_r(x_0)\times B_s(y_0)$ the two descriptions coincide exactly: the
zero set of $f$ and the graph of $g$ are the same set. Outside the box the level set may do
anything at all, including doubling back or carrying further branches, and that is precisely
why the statement is local.

\begin{center}
\begin{tikzpicture}[>=Latex, scale=1.2]
  \draw[green!60!black, thick, ->] (-0.4,0) -- (5.3,0) node[right, font=\small] {$x \in \mathbb{R}^d$};
  \draw[orange!90!black, thick, ->] (0,-0.4) -- (0,3.4) node[above, font=\small] {$y \in \mathbb{R}^{n-d}$};

  % Level set: inside the box a clean graph; outside it doubles back and branches.
  \draw[gray!55, very thick, smooth]
    plot coordinates {(0.45,2.95) (0.8,2.35) (1.1,1.98) (1.4,1.751)};
  \draw[gray!55, very thick, smooth]
    plot coordinates {(3.6,1.531) (4.05,1.33) (4.35,0.78) (4.15,0.28)};
  \node[gray!55!black, font=\scriptsize, right] at (4.45,0.80) {further branches};

  % The box, then the graph redrawn on top so it reads as the box's content.
  \draw[green!60!black, thick, fill=green!7] (1.4,0.9) rectangle (3.6,2.4);
  \draw[blue!80!black, very thick, domain=1.4:3.6, samples=60, smooth]
    plot ({\x}, {1.72 - 0.30*(\x-1.5) + 0.10*(\x-1.5)*(\x-1.5)});

  \coordinate (P) at (2.5,1.52);
  \fill (P) circle (1.7pt) node[above left, font=\scriptsize] {$(x_0,y_0)$};

  % Box extents.
  \draw[green!60!black, <->] (1.4,0.62) -- node[below, font=\scriptsize] {$B_r(x_0)$} (3.6,0.62);
  \draw[orange!90!black, <->] (1.13,0.9) -- node[left, font=\scriptsize] {$B_s(y_0)$} (1.13,2.4);
  \draw[dotted, gray] (2.5,0) -- (P);
  \node[below, font=\scriptsize] at (2.5,-0.05) {$x_0$};

  \node[blue!80!black, font=\scriptsize, right, align=left] at (3.72,2.55)
    {inside the box:\\ $f(x,y)=0 \iff y=g(x)$};
  \draw[blue!80!black, thin, ->] (3.72,2.35) -- (3.05,1.62);
\end{tikzpicture}
\end{center}

\begin{example}[The circle $S^1$]
The circle $S^1 = \{(x,y) \in \mathbb{R}^2 : x^2+y^2-1=0\}$ is the level set at $0$ of
$F(x,y) := x^2+y^2-1$, i.e., $S^1 = F^{-1}\{0\}$. Locally, there are explicit functions
\[x \mapsto \pm\sqrt{1-x^2} = y(x),\]
defined on open intervals around $x_0$ whenever $y(x_0)\neq0$. And indeed,
\[D_yF(x,y) = \frac{\partial F}{\partial y} = 2y\]
is bijective whenever $y \neq 0$, so the implicit function theorem holds there.

\begin{remark}[The circle realizes both panels]
The circle exhibits both cases of the figure above, and it is worth checking them against each
other.

At a point with $y_0 \neq 0$ --- say $(0,1)$ --- we have $D_yF = 2y_0 \neq 0$, the circle meets
that point at a non-vertical angle, and near it the circle really is the graph
$y = \sqrt{1-x^2}$. That is the left panel.

At $\bar x = 1$, where $y = 0$, we get $D_yF(1,0) = 0$. The circle is \emph{vertical} there, and
no graph $y(x)$ can exist around it: for $x$ slightly less than $1$ there are two points of the
circle above $x$, namely $\pm\sqrt{1-x^2}$, and for $x$ slightly greater than $1$ there are none.
That is the right panel exactly --- two $y$'s on one side, none on the other.

Note what has \emph{not} gone wrong at $\bar x = 1$: the circle is a perfectly smooth curve
there, with no corner and no break. What fails is only the attempt to describe it as a graph
\emph{over $x$}. Swap the roles of the two variables and everything works again, since
$D_xF(1,0) = 2 \neq 0$: near $(1,0)$ the circle is the graph $x = \sqrt{1-y^2}$ of a function
of $y$. This is the \qt{which variables get solved for is a choice} point above, in its
smallest possible instance.
\end{remark}
\end{example}

% Sonnet 5 (Medium) --- FIG-W06-06, drawn from the page image seen while transcribing
\begin{center}
\begin{tikzpicture}[>=Latex, scale=1.3]
  \draw[->] (-1.6,0) -- (1.9,0) node[right, font=\small] {$x$};
  \draw[->] (0,-1.6) -- (0,1.6) node[above, font=\small] {$y$};
  \draw[green!60!black, thick] (0,0) circle (1);

  \draw[red!80!black, dashed, thick, domain=-0.9:0.9, samples=40, smooth]
    plot ({\x}, {sqrt(1-\x*\x)}) node[right, font=\scriptsize] {$y=\sqrt{1-x^2}$};

  \coordinate (X0) at (0.4,0);
  \draw[dotted, gray] (X0) -- (0.4,0.917);
  \node[below, font=\scriptsize] at (X0) {$x_0$};

  \fill[red!80!black] (1,0) circle (1.5pt) node[below right, font=\scriptsize] {$\bar x$};
  \node[red!80!black, font=\scriptsize, below right, align=left] at (1.0,-0.5)
    {not possible\\around $y=0$};

  \draw[->, thick] (2.6,-1.0) -- (2.6,1.3) node[above, font=\small] {$\mathbb{R}$};
  \fill (2.6,0) circle (1.5pt) node[right, font=\small] {$0$};
  \draw[->, thick, blue!80!black] (1.15,0.6).. controls (2.0,1.0).. (2.6,0)
    node[midway, above, font=\small] {$F$};
\end{tikzpicture}
\end{center}

\begin{example}[An implicit quintic]
Let $G(x,y) := y^5+y^3+y+x$, and consider its level set at $0$,
$G^{-1}\{0\} = \{(x,y)\in\mathbb{R}^2 : G(x,y)=0\}$. Clearly $(0,0)\in G^{-1}\{0\}$, since
$G(0,0)=0^5+0^3+0+0=0$. Now both
\[\partial_xG(0,0)=1\neq0, \qquad \partial_yG(0,0)=\big[5y^4+3y^2+1\big]_{y=0}=1\neq0,\]
so by the implicit function theorem there are open intervals (balls) $(-r,r)$, $(-s,s)$ and
functions
\[(-r,r)\to(-s,s), \quad y\mapsto x(y) = -y^5-y^3-y,\]
\[(-s,s)\to(-r,r), \quad x\mapsto y(x),\]
satisfying
\[y^5+y^3+y+x(y)=0, \qquad y(x)^5+y(x)^3+y(x)+x=0.\]

\begin{remark}
As Galois and Abel showed, no closed-form expression for $y(x)$ in radicals exists in general,
since $y^5+y^3+y+x=0$ is a genuine quintic in $y$. Nevertheless, the implicit function theorem
guarantees that $y(x)$ exists (and is $C^\infty$, since $G$ is a polynomial), and it can be
computed and plotted numerically.
\end{remark}
\end{example}

% Sonnet 5 (Medium) --- FIG-W06-07, drawn from the page image seen while transcribing
\begin{center}
\begin{tikzpicture}[>=Latex, scale=0.62]
  \draw[->] (-11,0) -- (11,0) node[right, font=\small] {$x$};
  \draw[->] (0,-2.0) -- (0,2.0) node[above, font=\small] {$y(x)$};
  \foreach \x in {-10,-7.5,-5,-2.5,2.5,5,7.5,10}
    \draw (\x,0.08) -- (\x,-0.08) node[below, font=\tiny] {$\x$};
  \foreach \y in {-1.5,-1.0,-0.5,0.5,1.0,1.5}
    \draw (0.15,\y) -- (-0.15,\y) node[left, font=\tiny] {$\y$};

  \draw[blue!80!black, thick, variable=\t, domain=-1.42:1.42, samples=120, smooth]
    plot ({-\t*\t*\t*\t*\t - \t*\t*\t - \t}, {\t});
\end{tikzpicture}
\end{center}


% --- exercise relocated here from the dissolved sheet 7 ---

% Source: Jérôme Paschoud/Notizen/Woche 7.1 Satz über implizite Funktionen.pdf, p. 5
% Generator: Gemini 3.6 Flash (High)
\begin{exercise}[Implicitly defined functions from a system]
\label{ex:implicit_system_yz_of_x}
Show that one can solve the system of equations
\[
\begin{cases}
-2x^2 + y^2 + z^2 = 0 \\
x^2 + e^{y-1} - 2y = 0
\end{cases}
\]
for $y, z$ as functions of $x$ (i.e.\ $y(x), z(x)$) close enough to the point $(1,1,1)$. Calculate $\begin{pmatrix} y'(1) \\ z'(1) \end{pmatrix}$. (cf.\ Michaels Bsp 8.3.5)
\end{exercise}

\begin{exercisesolution}[\cref{ex:implicit_system_yz_of_x}]
Let $F(x,y,z) = \begin{pmatrix} -2x^2 + y^2 + z^2 \\ x^2 + e^{y-1} - 2y \end{pmatrix}$. We have $F(1,1,1) = \begin{pmatrix} -2+1+1 \\ 1+1-2 \end{pmatrix} = \begin{pmatrix} 0 \\ 0 \end{pmatrix}$.
The Jacobian of $F$ with respect to $(y,z)$ is
\[
D_{y,z} F(x,y,z) = \begin{pmatrix} \partial_y F_1 & \partial_z F_1 \\ \partial_y F_2 & \partial_z F_2 \end{pmatrix} = \begin{pmatrix} 2y & 2z \\ e^{y-1} - 2 & 0 \end{pmatrix}.
\]
At $(x,y,z) = (1,1,1)$, we have
\[
D_{y,z} F(1,1,1) = \begin{pmatrix} 2 & 2 \\ e^0 - 2 & 0 \end{pmatrix} = \begin{pmatrix} 2 & 2 \\ -1 & 0 \end{pmatrix}.
\]
The determinant is $\det(D_{y,z} F(1,1,1)) = 2(0) - 2(-1) = 2 \neq 0$. Thus the matrix is invertible, and by the Implicit Function Theorem, there exist functions $y(x), z(x)$ solving the system in a neighbourhood of $x=1$ such that $y(1)=1$ and $z(1)=1$.

To find the derivatives at $x=1$, we use the formula
\[
\begin{pmatrix} y'(1) \\ z'(1) \end{pmatrix} = - (D_{y,z} F(1,1,1))^{-1} D_x F(1,1,1).
\]
We calculate $D_x F(x,y,z) = \begin{pmatrix} -4x \\ 2x \end{pmatrix}$, so $D_x F(1,1,1) = \begin{pmatrix} -4 \\ 2 \end{pmatrix}$.
The inverse of $D_{y,z} F(1,1,1)$ is
\[
\begin{pmatrix} 2 & 2 \\ -1 & 0 \end{pmatrix}^{-1} = \frac{1}{2} \begin{pmatrix} 0 & -2 \\ 1 & 2 \end{pmatrix}.
\]
Then
\[
\begin{pmatrix} y'(1) \\ z'(1) \end{pmatrix} = - \frac{1}{2} \begin{pmatrix} 0 & -2 \\ 1 & 2 \end{pmatrix} \begin{pmatrix} -4 \\ 2 \end{pmatrix} = - \frac{1}{2} \begin{pmatrix} -4 \\ 0 \end{pmatrix} = \begin{pmatrix} 2 \\ 0 \end{pmatrix}.
\]
Therefore, $y'(1) = 2$ and $z'(1) = 0$.
\end{exercisesolution}

% Source: Jérôme Paschoud/Notizen/Woche 7.1 Satz über implizite Funktionen.pdf, p. 6
% Generator: Gemini 3.6 Flash (High)
\begin{exercise}[Implicit solving in different variables]
\label{ex:implicit_xyz_equation}
Consider the equation $xyz - yz^4 + x^2y = 2$ near the point $(0,-2,1)$. Can one solve this equation for:
\begin{enumerate}[label=\textbf{(\alph*)}]
  \item $x$?
  \item $y$?
  \item $z$?
\end{enumerate}
Calculate the derivative of the solution at the given point, if possible. (cf.\ Michaels Bsp 8.3.3)
\end{exercise}

\begin{exercisesolution}[\cref{ex:implicit_xyz_equation}]
Let $F(x,y,z) = xyz - yz^4 + x^2y - 2$. At the point $P = (0,-2,1)$, we have $F(0,-2,1) = 0 - (-2)(1) + 0 - 2 = 0$.
We compute the partial derivatives of $F$:
\begin{align*}
\partial_x F(x,y,z) &= yz + 2xy \\
\partial_y F(x,y,z) &= xz - z^4 + x^2 \\
\partial_z F(x,y,z) &= xy - 4yz^3
\end{align*}
Evaluating at $P = (0,-2,1)$:
\begin{align*}
\partial_x F(0,-2,1) &= (-2)(1) + 0 = -2 \neq 0 \\
\partial_y F(0,-2,1) &= 0 - 1 + 0 = -1 \neq 0 \\
\partial_z F(0,-2,1) &= 0 - 4(-2)(1) = 8 \neq 0
\end{align*}

\textbf{(a)} Since $\partial_x F(P) = -2 \neq 0$, the Implicit Function Theorem guarantees we can solve for $x = x(y,z)$ locally. The derivative (gradient) is
\[
\nabla x(y,z) \Big|_{(-2,1)} = \begin{pmatrix} \partial_y x \\ \partial_z x \end{pmatrix} = -\frac{1}{\partial_x F(P)} \begin{pmatrix} \partial_y F(P) \\ \partial_z F(P) \end{pmatrix} = -\frac{1}{-2} \begin{pmatrix} -1 \\ 8 \end{pmatrix} = \begin{pmatrix} -1/2 \\ 4 \end{pmatrix}.
\]

\textbf{(b)} Since $\partial_y F(P) = -1 \neq 0$, we can solve for $y = y(x,z)$ locally. The derivative is
\[
\nabla y(x,z) \Big|_{(0,1)} = \begin{pmatrix} \partial_x y \\ \partial_z y \end{pmatrix} = -\frac{1}{\partial_y F(P)} \begin{pmatrix} \partial_x F(P) \\ \partial_z F(P) \end{pmatrix} = -\frac{1}{-1} \begin{pmatrix} -2 \\ 8 \end{pmatrix} = \begin{pmatrix} -2 \\ 8 \end{pmatrix}.
\]

\textbf{(c)} Since $\partial_z F(P) = 8 \neq 0$, we can solve for $z = z(x,y)$ locally. The derivative is
\[
\nabla z(x,y) \Big|_{(0,-2)} = \begin{pmatrix} \partial_x z \\ \partial_y z \end{pmatrix} = -\frac{1}{\partial_z F(P)} \begin{pmatrix} \partial_x F(P) \\ \partial_y F(P) \end{pmatrix} = -\frac{1}{8} \begin{pmatrix} -2 \\ -1 \end{pmatrix} = \begin{pmatrix} 1/4 \\ 1/8 \end{pmatrix}.
\]
\end{exercisesolution}

% Originally: content/week-07.tex, \exercisesheet{7}, problem 7.5
% Source: exercises/Ex7_Analysis2_eng.pdf

\begin{exercise}[7.5 --- Implicit function II \textnormal{(important)}]
\label{ex:7.5}
Show that the system of equations
\[\begin{cases} xy^5+yu^5+zv^5=1, \\ x^5y+y^5u+z^5v=1, \end{cases}\]
is solvable for the variables $u$ and $v$ in a neighborhood of the point
$(x_0,y_0,z_0,u_0,v_0) = (0,1,1,1,0)$ and determine the derivatives $D_{(0,1,1)}u$ and
$D_{(0,1,1)}v$ of the implicitly defined functions $u=u(x,y,z)$ and $v=v(x,y,z)$.
\end{exercise}
